\documentclass[10pt,a4paper]{article}

\usepackage[T1]{fontenc}
\usepackage[utf8]{inputenc}
\usepackage{lmodern}
\usepackage[margin=1.6cm]{geometry}
\usepackage{xcolor}
\usepackage{array}
\usepackage{booktabs}
\usepackage{colortbl}
\usepackage{makecell}
\usepackage{enumitem}

% ---------- colours ----------
\definecolor{ink}{HTML}{1F2937}
\definecolor{accent}{HTML}{1D4ED8}
\definecolor{muted}{HTML}{6B7280}
\definecolor{headbg}{HTML}{F3F4F6}
\definecolor{greenrow}{HTML}{DCFCE7}
\definecolor{rulegray}{HTML}{D1D5DB}

\AtBeginDocument{\color{ink}}
\setlength{\parindent}{0pt}
\setlength{\parskip}{0.3em}
\renewcommand{\arraystretch}{1.1}
\arrayrulecolor{rulegray}

% ---------- helpers ----------
\newcommand{\code}[1]{\texttt{#1}}
\newcommand{\ind}{\hspace*{1.2em}}
\newcommand{\head}[1]{\vspace{0.6em}{\large\bfseries\color{accent}#1}\par\vspace{0.2em}}
\newcommand{\note}[1]{{\small\color{muted}#1}}
\newcolumntype{L}[1]{>{\raggedright\arraybackslash}p{#1}}

\begin{document}
\pagestyle{empty}

{\LARGE\bfseries Structure cleanup: the 5 rules}\par
\vspace{0.2em}
{\color{muted}\code{optimize\_prog\_structure} in \code{optimizer\_simple\_egglog.py},
lines 917--947 (branch dead-code-elimination). Every example below was run
through the real function.}

% ------------------------------------------------------------
\head{What it does}

A purely syntactic pass over the Prog tree. It works bottom-up: the body of
every \textbf{Cond} and every element of a \textbf{Sequence} is cleaned first,
then the rules are applied to the cleaned result. It looks at no variable
values. The pipeline runs it three times (after the first dead-code pass,
after the ite lowering, after the second dead-code pass, lines 1939, 1941,
1943) and once per round inside the dead-code loop (line 1909).

% ------------------------------------------------------------
\head{The rules}

\begingroup
\small
\noindent
\begin{tabular}{@{}l L{3.3cm} L{3.7cm} L{3.2cm} L{3.1cm}@{}}
\rowcolor{headbg}
\textbf{\#} & \textbf{Rule} & \textbf{Before} & \textbf{After} & \textbf{Why it is safe} \\
\midrule
1 & A guard that is the constant \textbf{true} is replaced by its body
  & \makecell[l]{\code{if (true) \{}\\\code{\ind y = 1;}\\\code{\}}}
  & \code{y = 1;}
  & The branch is always taken. \\
\midrule
2 & A guard that is the constant \textbf{false} removes the whole \textbf{if}
  & \makecell[l]{\code{if (false) \{}\\\code{\ind y = 1;}\\\code{\}}}
  & (nothing)
  & The branch is never taken. \\
\midrule
3 & An \textbf{if} whose (cleaned) body is empty is removed
  & \makecell[l]{\code{if (x < 0) \{}\\\code{\}}}
  & (nothing)
  & Evaluating a guard has no effect: terms are side-effect free. \\
\midrule
4 & A block inside a block is \textbf{flattened}
  & \makecell[l]{\code{\{ a = 1;}\\\code{\ind \{ b = 2; c = 3; \}}\\\code{\ind d = 4; \}}}
  & \makecell[l]{\code{\{ a = 1;}\\\code{\ind b = 2; c = 3;}\\\code{\ind d = 4; \}}}
  & Sequencing is associative. \\
\midrule
5 & An \textbf{empty block} inside a block is dropped
  & \makecell[l]{\code{\{ a = 1;}\\\code{\ind \{ \}}\\\code{\ind b = 2; \}}}
  & \makecell[l]{\code{\{ a = 1;}\\\code{\ind b = 2; \}}}
  & The empty block does nothing (it is the unit of sequencing). \\
\bottomrule
\end{tabular}
\endgroup

\vspace{0.3em}
\note{Rule 5 is the same effect as rule 4: flattening an empty block already
adds nothing to the list. The extra filter that drops empty blocks in the code
never changed a result in 20,000 random programs.}

% ------------------------------------------------------------
\head{Worked example (the innermost part is cleaned first)}

The inner guard becomes empty, so the guard above it becomes empty, and so on
up the tree. The result is still a block with one statement: a one-element
block is not unwrapped.

\vspace{0.4em}
\begingroup
\small
\noindent
\begin{tabular}{@{}l l l@{}}
\rowcolor{headbg}
\textbf{Step} & \textbf{Rule} & \textbf{Program becomes} \\
\midrule
start & & \code{if (true) \{ if (x < 0) \{ if (false) \{ y = 1; \} \} z = 2; \}} \\
\midrule
1 & 2 & \code{if (true) \{ if (x < 0) \{ \} z = 2; \}} \\
\midrule
2 & 3 & \code{if (true) \{ \{ \} z = 2; \}} \\
\midrule
3 & 4 / 5 & \code{if (true) \{ z = 2; \}} \\
\midrule
\rowcolor{greenrow}
4 & 1 & \code{\{ z = 2; \}} \\
\bottomrule
\end{tabular}
\endgroup

% ------------------------------------------------------------
\head{What it leaves alone}

\begin{itemize}[leftmargin=1.2em,itemsep=0.2em,topsep=0.2em,label=\textbullet]
  \item \textbf{Empty infinite loops} stay: \code{loop \{ \}} never reads and
        never ends, so removing it would change the behaviour.
  \item \textbf{One-element blocks} are not unwrapped, and an assignment with
        an empty list is not removed.
  \item \textbf{Non-constant guards} are never decided here. Terms like
        \code{x < x} are simplified by Egglog (rule \#75), and guards decided
        by known values are handled by the dead-code rule R2.
\end{itemize}

% ------------------------------------------------------------
\head{Checks and meaning}

\begin{itemize}[leftmargin=1.2em,itemsep=0.2em,topsep=0.2em,label=\textbullet]
  \item Running the pass twice gives the same result as running it once
        (5,000 random programs).
  \item All five rules change only \textbf{internal steps}, never what the
        environment sees at a Read and never how the program ends: in the
        model of \code{Prog.v} they are equivalences up to internal steps
        (weak bisimulation). That statement is argued here, \textbf{not yet
        proved in Coq}; the Coq proof so far covers only the Egglog stage.
\end{itemize}

\end{document}
